\documentclass{article}
%\usepackage{times}
%\usepackage{mathtime} 
\usepackage{pslatex} %regular times font/smaller courier font
\usepackage{amssymb}
\usepackage{alltt}
\usepackage{longtable}
\usepackage{multirow}
\usepackage{url}
\usepackage{html}

%\newtheorem{thm}{Theorem}[section]
\newtheorem{cor}{Corollary}
\newtheorem{lem}{Lemma}
\newtheorem{defin}{Definition}
%Use like this: \begin{lem}\label{...}  .... \end{lem}

\renewcommand{\labelitemi}{ } %to make itemize have bullets instead of confusing -'s
\pagestyle{plain} %just for editing purposes

\renewcommand{\phi}{\varphi}
\newcommand{\always}{\Box}
\newcommand{\eventually}{\Diamond}
\newcommand{\calL}{{\cal L}}


%%%%%%%%%%%%%%%%%%%%%%%%%%%%%%%%%%%%%%%%%%%%%%%%%%%%%%%%
% Figure Magic
%%%%%%%%%%%%%%%%%%%%%%%%%%%%%%%%%%%%%%%%%%%%%%%%%%%%%%%%
\usepackage{epsfig}
\usepackage{float}
\usepackage{subfigure}
\renewcommand{\topfraction}{.95} %figures can take up at most 95% of the page before being alone
\renewcommand{\bottomfraction}{.99} %figures can take up at most 99% of the page before being alone
\renewcommand{\textfraction}{.1} %at most this this % of page will be text before making figure-only page
\addtolength{\textfloatsep}{-5mm}
\addtolength{\floatsep}{-5mm}
\addtolength{\intextsep}{-5mm}
\addtolength{\abovecaptionskip}{-5mm}
\addtolength{\belowcaptionskip}{-3mm}

\title{Model Checking Benchmarking Scripts}
\author{Kristin Y. Rozier}
\date{2007}

\begin{document}

\maketitle

These pages contain further details of the experiments described in \htmladdnormallink{``LTL Satisfiability Checking''}{http://shemesh.larc.nasa.gov/people/kyr/publications.html} by K.Y.Rozier and M.Y.Vardi.

All of the software available for download on this page was written by Kristin Yvonne Rozier. It is released under a \htmladdnormallink{NASA Software Agreement}{http://shemesh.larc.nasa.gov/people/kyr/software_agreement.txt}. Please feel free to email me concerning clarifications, bugs, or other corrections.

\section{Explicit Tool Models}
To benchmark LTL formula translations, we use for the SPIN model a \emph{universal} Promela program, enumerating all possible traces over $Prop$. For example, when $Prop = \{a, b\}$, (and therefore the number $n$ of variables is 2), the Promela model is:

\begin{verbatim}
/*File: 2vars.pml*/

bool A,B;

/* define an active procedure to generate values for A and B */

active proctype generateValues()
{ do
   :: atomic{ A = 0; B = 0; }
   :: atomic{ A = 0; B = 1; }
   :: atomic{ A = 1; B = 0; }
   :: atomic{ A = 1; B = 1; }
  od }
\end{verbatim}

We use the \texttt{atomic\{\}} construct to ensure that the Boolean variables change value in one unbreakable step. Note that the size of this model is exponential in the number of atomic propositions.

The variables in $Prop$ are defined in the accompanying file \texttt{2vars.define}, which is prepended to the never-claim we are testing.

These two files are generated automatically for any given value of $n$ by any one of the following scripts. (Just for simplicity, there are 3 versions of the same script here.):

\medskip\noindent
\begin{tabular}{lll}
\multicolumn{1}{c}{Script Name} & \multicolumn{1}{c}{Variable Set} & \multicolumn{1}{c}{Recommended Input Formulas} \\
\htmladdnormallink{generateValues.pl}{generateValues.pl.zip} & Prop = {a, b, c, \ldots, aa, bb, cc, \ldots} & Use with random and counter formulas. \\
\htmladdnormallink{generatePvalues.pl}{generatePvalues.pl.zip} & Prop = {p1, p2, p3, \ldots} & Use with pattern formulas. \\
\htmladdnormallink{generateAlphaPvalues.pl}{generateAlphaPvalues.pl.zip} & Prop = {po, pt, ph, \ldots} & Use with pattern formulas, for tools \\
 & & which don't allow numbers in variable names \\
\end{tabular}

\medskip \noindent
Usage: \texttt{generateValues.pl number\_of\_variables}

\medskip \noindent
Here is an example of how to use these scripts to run SPIN on a generated Promela never-claim with a universal model:

\begin{alltt}
#Preceed the formula with a 'next time' operator to skip 
#  SPIN's nondeterministic assignment upon declaration
$LTL_to_automata_tool_binary ``(X($input_formula))" >& $temp #generate the automaton
generateValues.pl $n
cp -f ${n}vars.define pan.ltl
cat $temp >> pan.ltl
rm -f $temp
cp -f ${n}vars.pml pan_in
spin -a -X -N pan.ltl pan_in
gcc -w -o pan -D_POSIX_SOURCE -DMEMLIM=128 -DXUSAFE -DNOFAIR  pan.c
./pan -v -X -m10000 -w19  -a -c1
\end{alltt}


\section{Symbolic Tool Models}

We can also benchmark LTL formula modeling with CadenceSMV and NuSMV. To check whether a LTL formula $\phi$ is satisfiable, we model check $\neg\phi$ against a universal SMV model. For example, if $\phi = (X (a))$, we provide the following input to NuSMV:
\small
\begin{verbatim}
MODULE main
  VAR
    a : boolean;
    b : boolean;
    c : boolean;
  LTLSPEC !(X(a=1 ))
  FAIRNESS
    1
\end{verbatim}
\normalsize
(The ``FAIRNESS'' line varies with versions of NuSMV. The CadenceSMV model is very similar.) SMV negates the specification, $\neg \phi$, symbolically compiles $\phi$ into $A_{\phi}$, and conjoins $A_{\phi}$ with the universal model. If the automaton is not empty, then SMV finds a fair path, which satisfies the formula $\phi$. In this way, SMV acts as both a symbolic compiler and a search engine.

\section{Input Formulas}
In \cite{RV07}, we benchmarked the tools against three types of scalable formulas: random formulas, counter formulas, and pattern formulas. Scalability played an important role in our experiment, since the goal was to challenge the tools with large formulas and state spaces.

\subsection{Random Formulas}

In order to cover as much of the problem space as possible, I wrote a script to generate sets of 500 randomly-generated formulas varying the formula length and number of variables as in \cite{DGV99}. It randomly generates sets of 500 formulas varying the number of variables, $N$, from 1 to 5, and the length of the formula, $L$, from 5 up to 200. It sets the probability of choosing a temporal operator $P = 0.5$ to create formulas with both a nontrivial temporal structure and a nontrivial Boolean structure. For temporal tests, formulas are also generated with $P = \frac{1}{3}$, $P = 0.7$, and $P = 0.95$. Other choices are decided uniformly. When generating a set of length $L$, every formula created is exactly of length $L$ and not \emph{up to} $L$, however, these randomly-generated formulas are frequently reducible.

\medskip \noindent
Here is the script used to generate random formulas in the style of \cite{DGV99}:
\begin{alltt}
\htmladdnormallink{generateFormulas.pl}{generateFormulas.pl.zip}
        With no arguments: generate a wide range of formula files
        With 3 arguments: generate files for specified L, N, P (order matters)
\end{alltt}

\medskip \noindent
In our paper, \cite{RV07}, all of the random formulas we used were generated prior to testing, so each tool was run on the \emph{same} formulas. If you wish to re-create our tests exactly, the set of formulas generated for our experiments are here:

\medskip \noindent
\htmladdnormallink{\texttt{formulas/}}{formulas/}\\
\htmladdnormallink{formulas.zip}{formulas.zip}


\subsection{Counter Formulas}

Pre-translation rewriting is highly effective for random formulas, but ineffective for structured formulas \cite{EH00,SB00}. To measure performance on scalable, non-random formulas we tested the tools on formulas that describe $n$-bit binary counters with increasing values of $n$. These formulas are irreducible by pre-translation rewriting, uniquely satisfiable, and represent a predictably-sized state space. Correctness is easily determined by checking for precisely the unique counterexample for each counter formula. We tested four constructions of binary counter formulas, varying two factors: number of variables and nesting of $\mathcal{X}$'s.

We can represent a binary counter using two variables: a counter variable and a marker variable to designate the beginning of each new counter value. Alternatively, we can use 3 variables, adding a variable to encode carry bits, which eliminates the need for $\mathcal{U}$-connectives in the formula. We can nest $\mathcal{X}$'s to provide more succinct formulas or express the formulas using a conjunction of unnested $\mathcal{X}$-sub-formulas.

Let $b$ be an atomic proposition. Then a computation $\pi$ over $b$ is a word in $(2^{\{0,1\}})^\omega$. By dividing $\pi$ into blocks of length $n$, we can view $\pi$ as a sequence of $n$-bit values, denoting the sequence of values assumed by an $n$-bit counter starting at $0$, and incrementing successively by $1$. To simplify the formulas, we represent each block $b_0,b_1,\ldots,b_{n-1}$ as having the most significant bit on the right and the least significant bit on the left. For example, for $n=2$ the $b$ blocks cycle through the values $00$, $10$, $01$, and $11$. For technical convenience, we use an atomic proposition $m$ to mark the blocks. That is, we intend $m$ to hold at point $i$ precisely when $i=0 \bmod n$.

For $\pi$ to represent an $n$-bit counter, the following properties need to hold:

\begin{alltt}
1) The marker consists of a repeated pattern of a 1 followed by n-1 0's.
2) The first n bits are 0's.
3) If the least significant bit is 0, then it is 1 n steps later 
   and the other bits do not change.
4) All of the bits before and including the first 0 in an n-bit block flip
   their values in the next block; the other bits do not change.
\end{alltt} 

For $n=4$, these properties are captured by the conjunction of the following formulas:

\begin{alltt}
1. (m) && ( [](m -> ((X(!m)) && (X(X(!m))) && (X(X(X(!m)))) 
                                           && X(X(X(X(m)))))))
2. (!b) && (X(!b)) && (X(X(!b))) && (X(X(X(!b))))
3. []( (m && !b) -> 
       ( X(X(X(X(b)))) && 
         X ( ( (!m) && 
           (b -> X(X(X(X(b))))) && 
           (!b -> X(X(X(X(!b))))) ) U m ) ) )
4. [] ( (m && b) -> 
        (X(X(X(X(!b)))) && 
          (X ( (b && !m && X(X(X(X(!b))))) U 
               (m || (!m && !b && X(X(X(X(b)))) && 
                      X( ( !m && (b -> X(X(X(X(b))))) && 
                          (!b -> X(X(X(X(!b))))) ) U m ) ) ) ) ) ) )
\end{alltt}

Scalable counter formulas of this form are automatically generated by \htmladdnormallink{LTLcounter.pl}{LTLcounter.pl.zip}. 
Usage: \texttt{LTLcounter.pl number\_of\_bits}



Note that this encoding creates formulas of length $O(n^2)$.
A more compact encoding results in formulas of length $O(n)$.
For example, we can replace formula (2) above with:

\small 
\begin{alltt}
2. ((!b) && X((!b) && X((!b) && X(!b))))
\end{alltt}
\normalsize

Scalable counter formulas of this form, with nested $\mathcal{X}$-operators, are automatically generated by \htmladdnormallink{LTLcounterLinear.pl}{LTLcounterLinear.pl.zip}. Usage: \texttt{LTLcounterLinear.pl number\_of\_bits}


We can eliminate the use of $\mathcal{U}$-connectives in the formula by adding an atomic proposition $c$ representing the carry bit. The required properties of an $n$-bit counter with carry are as follows:
\begin{center}
\begin{tabular}{c | cccccccccc}
\multicolumn{11}{c}{\bf A 4-bit Binary Counter} \\
\hline
m & 0001 & 0001 & 0001 & 0001 & 0001 & 0001 & 0001 & 0001 & 0001 & 0001 \\
b & 0000 & 0001 & 0010 & 0011 & 0100 & 0101 & 0110 & 0111 & 1000 & 1001 \\
c & 0000 & 0001 & 0000 & 0011 & 0000 & 0001 & 0000 & 0111 & 0000 & 0001 \\
\hline
\multicolumn{11}{c}{ }\\
\hline
m & 0001 & 0001 & 0001 & 0001 & 0001 & 0001 & 0001 & \ldots \\
b & 1010 & 1011 & 1100 & 1101 & 1110 & 1111 & 0000 & \ldots \\
c & 0000 & 0011 & 0000 & 0001 & 0000 & 1111 & 0000 & \ldots \\
\hline
\end{tabular}
\end{center}

\medskip \noindent
To describe the behavior of the counter, we describe six properties and \texttt{AND} them together. For example, here is the program output for a 4-bit counter:

\begin{alltt}
1) The marker consists of a repeated pattern of a 1 followed by n-1 0's.
2) The first n bits are 0's.
3) If m is 1 and b is 0 then c is 0 and n steps later b is 1.
4) If m is 1 and b is 1 then c is 1 and n steps later b is 0.
5) If there is no carry, then the next bit stays the same n steps later.
6) If there is a carry, flip the next bit n steps later and adjust the carry.
\end{alltt}

For $n=4$, these properties are captured by the conjunction of the following formulas.

\begin{alltt}
1. (m) && ( [](m -> ((X(!m)) && (X(X(!m))) && (X(X(X(!m)))) 
                                           && (X(X(X(X(m))))))))
2. (!b) && (X(!b)) && (X(X(!b))) && (X(X(X(!b))))
3. [] ( (m && !b) -> (!c && X(X(X(X(b))))) )
4. [] ( (m && b) -> (c && X(X(X(X(!b))))) )
5. [] (!c & X(!m)) -> 
      ( X(!c) && (X(b) -> X(X(X(X(b))))) && 
                 (X(!b) -> X(X(X(X(!b))))) )
6. [] (c -> ( ( X(!b) -> ( X(!c) && X(X(X(X(!b)))) ) ) && 
                         ( X(c) && X(X(X(X(b)))) ) ))
\end{alltt}

Scalable counter formulas of this form, using a carry bit and no $\mathcal{U}$-operators, are automatically generated by \htmladdnormallink{LTLcounterCarry.pl}{LTLcounterCarry.pl.zip}. Usage: \texttt{LTLcounterCarry.pl number\_of\_bits}


Scalable counter formulas using a carry bit and nested $\mathcal{X}$-operators are automatically generated by \htmladdnormallink{LTLcounterCarryLinear.pl}{LTLcounterCarryLinear.pl.zip}. Usage: \texttt{LTLcounterCarryLinear.pl number\_of\_bits}



\subsection{Pattern Formulas}

These scripts generate formulas from the eight classes of scalable formulas defined by \cite{GH06} to evaluate the performance of model checking algorithms on temporally-complex formulas. 

\medskip \noindent
Run them like this:\\
\texttt{formula.pl [-v] n}\\
where $n$ is the number of variables (ie the size) of the formula you desire. The flag -v is for ``verbose.''

\[E(n) = \bigwedge_{i=1}^{n} \Diamond p_i \]
\[U(n) = (\ldots (p_1 \ \mathcal{U} \ p_2) \ \mathcal{U} \ \ldots) \ \mathcal{U} \ p_n \]
\[R(n) = \bigwedge_{i=1}^{n} (\Box \Diamond p_i \vee \Diamond \Box p_{i+1}) \]
\[U_2(n) = p_1 \ \mathcal{U} \ (p_2 \ \mathcal{U} \ (\ldots p_{n-1} \ \mathcal{U} \ p_n) \ldots ) \]
\[C_1(n) = \bigvee_{i=1}^{n} \Box \Diamond p_i \]
\[C_2(n) = \bigwedge_{i=1}^{n} \Box \Diamond p_i \]
\[Q(n) = \bigwedge (\Diamond p_i \vee \Box p_{i+1}) \]
\[S(n) = \bigwedge_{i=1}^{n} \Box p_i \]

\begin{enumerate}
\item \htmladdnormallink{Eformula.pl}{Eformula.pl.zip}
\item \htmladdnormallink{Uformula.pl}{Uformula.pl.zip}
\item \htmladdnormallink{Rformula.pl}{Rformula.pl.zip}
\item \htmladdnormallink{U2formula.pl}{U2formula.pl.zip}
\item \htmladdnormallink{C1formula.pl}{C1formula.pl.zip}
\item \htmladdnormallink{C2formula.pl}{C2formula.pl.zip}
\item \htmladdnormallink{Qformula.pl}{Qformula.pl.zip}
\item \htmladdnormallink{Sformula.pl}{Sformula.pl.zip}
\end{enumerate}


\bibliographystyle{plain}
\bibliography{benchmarking_scripts.bib}

\end{document}
